% ID: MAT-TRI-EQG-014
% Titolo: Equazioni goniometriche e periodicità (1)
% Disciplina: Matematica
% Area: trigonometria
% Argomento: equazioni_goniometriche
% Classe: Quarta secondaria di secondo grado
% Difficolta: 3
% Tipo: problema_numerico
% Tag: matematica, trigonometria, equazioni_goniometriche
% Fonte: verifica_goniometria_base_anonimizzata
% Provenienza: materiale fornito dal docente; testo anonimizzato; soluzione redatta per Exergo

\begin{esercizio}
Risolvi le seguenti equazioni: 

\begin{enumerate}
\item $\displaystyle\tan\left(2x-\dfrac{\pi}{6}\right)
 =\tan\left(x+\dfrac{\pi}{4}\right);$

\item $\displaystyle 2\cos x+\sqrt{2}=0.$
\end{enumerate}
Salvo diversa indicazione, lavora in $\mathbb R$; angoli in radianti.
\end{esercizio}

\begin{soluzione}
a) Per l'uguaglianza delle tangenti $2x-\pi/6=x+\pi/4+k\pi$,
quindi $x=5\pi/12+k\pi$. Entrambe le tangenti sono definite per questi valori.
b) $\cos x=-\sqrt2/2$: $x=3\pi/4+2k\pi$ oppure $x=5\pi/4+2k\pi$. In tutte le famiglie di soluzioni $k\in\mathbb{Z}$.
\end{soluzione}
